\documentclass[aspectratio=169,12pt]{beamer}
% Estilo canónico para las presentaciones de Métodos Numéricos 2026-2.
\usepackage[T1]{fontenc}
\usepackage[utf8]{inputenc}
\usepackage[spanish,es-nodecimaldot]{babel}
\usepackage{lmodern}
\usepackage{amsmath,amssymb,mathtools}
\usepackage{booktabs}
\usepackage{graphicx}
\usepackage{xcolor}
\usepackage{hyperref}

\definecolor{upazul}{HTML}{003B70}
\definecolor{upazulclaro}{HTML}{EAF2F8}
\definecolor{updorado}{HTML}{C5962E}
\definecolor{upgris}{HTML}{4A5560}
\definecolor{upverde}{HTML}{167C5A}
\definecolor{uprojo}{HTML}{B33A3A}

\setbeamercolor{normal text}{fg=black,bg=white}
\setbeamercolor{structure}{fg=upazul}
\setbeamercolor{frametitle}{fg=white,bg=upazul}
\setbeamercolor{title}{fg=upazul}
\setbeamercolor{subtitle}{fg=upgris}
\setbeamercolor{block title}{fg=white,bg=upazul}
\setbeamercolor{block body}{fg=black,bg=upazulclaro}
\setbeamercolor{alerted text}{fg=uprojo}
\setbeamercolor{example text}{fg=upverde}

\setbeamerfont{title}{size=\Huge,series=\bfseries}
\setbeamerfont{subtitle}{size=\Large}
\setbeamerfont{frametitle}{size=\LARGE,series=\bfseries}
\setbeamerfont{normal text}{size=\large}
\setbeamerfont{block title}{series=\bfseries}

\setbeamertemplate{navigation symbols}{}
\setbeamertemplate{itemize item}{\color{updorado}$\blacktriangleright$}
\setbeamertemplate{itemize subitem}{\color{upazul}$\bullet$}
\setbeamertemplate{blocks}[rounded][shadow=false]

\setbeamertemplate{frametitle}{%
  \nointerlineskip
  \begin{beamercolorbox}[wd=\paperwidth,ht=0.78cm,dp=0.24cm,leftskip=0.55cm]{frametitle}
    \usebeamerfont{frametitle}\insertframetitle
  \end{beamercolorbox}}

\setbeamertemplate{footline}{%
  \leavevmode
  \hbox{%
    \begin{beamercolorbox}[wd=.82\paperwidth,ht=2.7ex,dp=1.2ex,leftskip=.5cm]{author in head/foot}
      \color{upgris}\scriptsize Métodos Numéricos 2026-2 \hspace{0.8em} | \hspace{0.8em} Universidad del Pacífico
    \end{beamercolorbox}%
    \begin{beamercolorbox}[wd=.18\paperwidth,ht=2.7ex,dp=1.2ex,rightskip=.5cm plus1fil]{date in head/foot}
      \color{upazul}\scriptsize\hfill\insertframenumber/\inserttotalframenumber
    \end{beamercolorbox}}
  \vskip0pt}

\newcommand{\concepto}[1]{\textcolor{upazul}{\textbf{#1}}}
\newcommand{\resultado}[1]{\textcolor{upverde}{\textbf{#1}}}
\newcommand{\advertencia}[1]{\textcolor{uprojo}{\textbf{#1}}}
\newcommand{\codigo}[1]{\colorbox{upazulclaro}{\strut\texttt{#1}}}
\newcommand{\fuente}[1]{\vfill\hfill{\tiny\color{upgris}#1}}



\usepackage{array}

\title{Sistemas de ecuaciones no lineales}
\subtitle{Punto fijo, Newton amortiguado y Broyden}
\author{Oswaldo José Velásquez Castañón}
\institute{Universidad del Pacífico}
\date{Métodos Numéricos -- 2026-2}

\newcommand{\R}{\mathbb{R}}
\newcommand{\mat}[1]{\begin{bmatrix}#1\end{bmatrix}}
\newcommand{\norm}[1]{\left\lVert#1\right\rVert}

\begin{document}

\begin{frame}[plain]
  \begin{columns}[T,onlytextwidth]
    \begin{column}{0.70\textwidth}
      \vspace{0.7cm}
      {\usebeamerfont{title}\color{upazul}\inserttitle\par}
      \vspace{0.35cm}
      {\usebeamerfont{subtitle}\color{upgris}\insertsubtitle\par}
      \vspace{1.1cm}
      {\large\insertauthor\par}
      {\normalsize\insertdate\par}
    \end{column}
    \begin{column}{0.25\textwidth}
      \vspace{0.55cm}
      \includegraphics[width=\linewidth]{UPlogo.pdf}
    \end{column}
  \end{columns}
\end{frame}

\begin{frame}{El problema vectorial}
Buscamos $x^*\in\R^n$ tal que
\[
F(x^*)=0,
\qquad
F=\mat{f_1\\ \vdots\\ f_n}:\R^n\to\R^n.
\]
En dos dimensiones, cada ecuación $f_i(x_1,x_2)=0$ describe una curva. Una
solución es un punto donde ambas curvas se intersectan.
\begin{block}{Dificultad adicional}
Las variables están acopladas. Cada actualización debe corregir todas las
ecuaciones al mismo tiempo.
\end{block}
\end{frame}

\begin{frame}{Punto fijo en varias variables}
Una reformulación $F(x)=0$ puede escribirse como
\[
x=G(x),
\qquad
x_{k+1}=G(x_k).
\]
Si $G$ lleva un conjunto cerrado $D$ dentro de sí mismo y
\[
\norm{G(x)-G(y)}\leq L\norm{x-y},
\qquad L<1,
\]
entonces la iteración converge al único punto fijo de $G$ en $D$.
\begin{block}{Condición local}
$\rho(J_G(x^*))<1$ garantiza convergencia para puntos iniciales cercanos.
\end{block}
\end{frame}

\begin{frame}{La norma elegida afecta la verificación}
Para una matriz $A$, la desigualdad
\[
\norm{G(x)-G(y)}\leq
\sup_{\xi\in D}\norm{J_G(\xi)}\norm{x-y}
\]
permite estimar una constante de Lipschitz.
\begin{itemize}
  \item Si alguna norma matricial da una cota menor que uno, tenemos una contracción.
  \item Una cota mayor que uno no demuestra divergencia.
  \item El radio espectral de $J_G(x^*)$ describe el comportamiento local.
\end{itemize}
Las normas son equivalentes en dimensión finita, pero sus cotas pueden tener
calidades muy distintas.
\end{frame}

\begin{frame}{Newton linealiza todas las ecuaciones}
La expansión de Taylor alrededor de $x_k$ es
\[
F(x_k+s)\approx F(x_k)+J_F(x_k)s.
\]
El paso de Newton anula el modelo lineal:
\[
J_F(x_k)s_k=-F(x_k),
\qquad
x_{k+1}=x_k+s_k.
\]
\begin{block}{Implementación correcta}
Se resuelve el sistema lineal para $s_k$. No se calcula $J_F(x_k)^{-1}$ de forma explícita.
\end{block}
\end{frame}

\begin{frame}{El jacobiano organiza las sensibilidades}
Para $F=(f_1,\ldots,f_n)^T$,
\[
J_F(x)=
\mat{
\dfrac{\partial f_1}{\partial x_1} & \cdots & \dfrac{\partial f_1}{\partial x_n}\\[2mm]
\vdots & \ddots & \vdots\\[1mm]
\dfrac{\partial f_n}{\partial x_1} & \cdots & \dfrac{\partial f_n}{\partial x_n}
}.
\]
La columna $j$ mide cómo cambian todas las ecuaciones al perturbar $x_j$.
\begin{block}{Lectura del paso}
$s_k$ compensa simultáneamente los residuos mediante el modelo local definido por $J_F(x_k)$.
\end{block}
\end{frame}

\begin{frame}{Algoritmo de Newton para sistemas}
\begin{block}{Entrada}
$F$, $J_F$, $x_0$, tolerancias y máximo de iteraciones.
\end{block}
Para $k=0,1,\ldots$:
\begin{enumerate}
  \item evaluar $r_k=F(x_k)$;
  \item resolver $J_F(x_k)s_k=-r_k$;
  \item actualizar $x_{k+1}=x_k+s_k$;
  \item comprobar $\norm{F(x_{k+1})}$ y $\norm{s_k}/\max\{1,\norm{x_{k+1}}\}$.
\end{enumerate}
Una factorización con pivoteo permite resolver el sistema lineal de cada paso.
\end{frame}

\begin{frame}{Ejemplo: circunferencia y recta}
Considere
\[
F(x,y)=
\mat{x^2+y^2-1\\x-y}=0.
\]
El jacobiano es
\[
J_F(x,y)=\mat{2x&2y\\1&-1}.
\]
Las soluciones son
\[
(x^*,y^*)=\left(\frac{1}{\sqrt2},\frac{1}{\sqrt2}\right),
\qquad
\left(-\frac{1}{\sqrt2},-\frac{1}{\sqrt2}\right).
\]
El punto inicial determina cuál de ellas atrae la iteración.
\end{frame}

\begin{frame}{Primer paso de Newton}
Desde $x_0=(1,0.5)^T$,
\[
F(x_0)=\mat{0.25\\0.5},
\qquad
J_F(x_0)=\mat{2&1\\1&-1}.
\]
El paso resuelve
\[
\mat{2&1\\1&-1}\mat{s_1\\s_2}
=-\mat{0.25\\0.5},
\qquad
s=\mat{-0.25\\0.25}.
\]
Por tanto,
\[
x_1=x_0+s=\mat{0.75\\0.75}.
\]
Una sola linealización elimina exactamente el residuo de la segunda ecuación.
\end{frame}

\begin{frame}{Iteraciones del ejemplo}
\centering
\begin{tabular}{cccc}
\toprule
$k$ & $x_k$ & $y_k$ & $\norm{F(x_k)}_2$ \\
\midrule
0 & 1.000000000 & 0.500000000 & $5.590\times10^{-1}$ \\
1 & 0.750000000 & 0.750000000 & $1.250\times10^{-1}$ \\
2 & 0.708333333 & 0.708333333 & $3.472\times10^{-3}$ \\
3 & 0.707107843 & 0.707107843 & $3.003\times10^{-6}$ \\
4 & 0.707106781 & 0.707106781 & $2.255\times10^{-12}$ \\
\bottomrule
\end{tabular}
\vspace{0.4cm}

La simetría $x=y$ aparece después del primer paso y reduce el resto del cálculo
a la aproximación de $1/\sqrt2$.
\end{frame}

\begin{frame}{Convergencia local de Newton}
Si $F(x^*)=0$, $F$ es continuamente diferenciable cerca de $x^*$ y
$J_F(x^*)$ es no singular, entonces Newton converge superlinealmente desde
puntos iniciales suficientemente cercanos.

Si además $J_F$ es Lipschitz cerca de $x^*$, entonces
\[
\norm{x_{k+1}-x^*}\leq C\norm{x_k-x^*}^2.
\]
\begin{block}{Alcance del resultado}
La rapidez está garantizada cerca de la solución. El teorema no ofrece una
garantía global para un punto inicial arbitrario.
\end{block}
\end{frame}

\begin{frame}{Jacobiano mal condicionado}
El paso satisface
\[
J_F(x_k)s_k=-F(x_k).
\]
Si $J_F(x_k)$ es casi singular:
\begin{itemize}
  \item una perturbación pequeña del residuo puede cambiar mucho el paso;
  \item el sistema lineal puede perder cifras significativas;
  \item la dirección calculada puede ser excesivamente grande.
\end{itemize}
\begin{block}{Diagnóstico}
Conviene vigilar el condicionamiento del jacobiano, el tamaño del paso y la
reducción efectiva del residuo.
\end{block}
\end{frame}

\begin{frame}{Newton amortiguado amplía la región útil}
Definimos la función de mérito
\[
\phi(x)=\frac12\norm{F(x)}_2^2.
\]
Después de calcular el paso de Newton $s_k$, usamos
\[
x_{k+1}=x_k+\alpha_k s_k,
\qquad 0<\alpha_k\leq1.
\]
La búsqueda hacia atrás reduce $\alpha_k$ hasta obtener una disminución
suficiente de $\phi$.
\begin{block}{Comportamiento esperado}
Lejos de la solución se amortigua el paso. Cerca de ella se acepta $\alpha_k=1$
y se recupera la velocidad de Newton.
\end{block}
\end{frame}

\begin{frame}{El jacobiano puede aproximarse por diferencias}
La columna $j$ puede aproximarse mediante
\[
J_F(x)e_j\approx
\frac{F(x+h_j e_j)-F(x)}{h_j}.
\]
Esto requiere aproximadamente $n+1$ evaluaciones de $F$ para construir un
jacobiano completo hacia adelante.
\begin{itemize}
  \item Un $h_j$ grande produce error de truncamiento.
  \item Un $h_j$ demasiado pequeño produce cancelación.
  \item La escala de cada componente debe influir en $h_j$.
\end{itemize}
La aproximación evita derivadas analíticas, pero todavía reconstruye una matriz en cada paso.
\end{frame}

\begin{frame}{Broyden generaliza la idea de la secante}
Sean
\[
s_k=x_{k+1}-x_k,
\qquad
y_k=F(x_{k+1})-F(x_k).
\]
El jacobiano exacto cumple aproximadamente $J_F(x_{k+1})s_k\approx y_k$.
Broyden construye matrices $B_k$ que satisfacen la condición secante
\[
B_{k+1}s_k=y_k.
\]
\begin{block}{Objetivo}
Actualizar una aproximación del jacobiano usando evaluaciones de $F$, sin
calcular todas sus derivadas en cada iteración.
\end{block}
\end{frame}

\begin{frame}{Actualización de Broyden}
La actualización de rango uno conocida como Broyden bueno es
\[
B_{k+1}
=B_k+
\frac{(y_k-B_ks_k)s_k^T}{s_k^Ts_k}.
\]
En efecto,
\[
B_{k+1}s_k
=B_ks_k+(y_k-B_ks_k)=y_k.
\]
La corrección modifica $B_k$ únicamente en la dirección necesaria para cumplir
la nueva condición secante.
\begin{block}{Costo}
Se evita reconstruir el jacobiano completo, aunque todavía debe resolverse un
sistema lineal con $B_k$.
\end{block}
\end{frame}

\begin{frame}{Algoritmo de Broyden}
Elegir $x_0$ y una matriz no singular $B_0$, por ejemplo $J_F(x_0)$.
Para $k=0,1,\ldots$:
\begin{enumerate}
  \item resolver $B_ks_k=-F(x_k)$;
  \item definir $x_{k+1}=x_k+s_k$;
  \item calcular $y_k=F(x_{k+1})-F(x_k)$;
  \item actualizar $B_{k+1}$ con la fórmula de rango uno;
  \item comprobar paso y residuo.
\end{enumerate}
Si $s_k^Ts_k$ es demasiado pequeño, se omite o reinicia la actualización.
\end{frame}

\begin{frame}{Broyden en el ejemplo de dos ecuaciones}
Con $x_0=(1,0.5)^T$ y $B_0=J_F(x_0)$:
\centering
\begin{tabular}{cccc}
\toprule
$k$ & $x_k$ & $y_k$ & $\norm{F(x_k)}_2$ \\
\midrule
0 & 1.000000000 & 0.500000000 & $5.590\times10^{-1}$ \\
1 & 0.750000000 & 0.750000000 & $1.250\times10^{-1}$ \\
2 & 0.708333333 & 0.708333333 & $3.472\times10^{-3}$ \\
3 & 0.707142857 & 0.707142857 & $1.020\times10^{-4}$ \\
4 & 0.707106812 & 0.707106812 & $8.842\times10^{-8}$ \\
5 & 0.707106781 & 0.707106781 & $2.255\times10^{-12}$ \\
\bottomrule
\end{tabular}

Broyden necesita más pasos que Newton, pero evita nuevas evaluaciones del jacobiano.
\end{frame}

\begin{frame}{Comparación de Newton y Broyden}
\centering
\small
\begin{tabular}{p{2.6cm}p{3.2cm}p{3.2cm}}
\toprule
Criterio & Newton & Broyden \\
\midrule
Derivadas & Jacobiano en cada paso & Matriz inicial y actualizaciones \\
Velocidad local & Cuadrática bajo hipótesis estándar & Habitualmente superlineal \\
Costo dominante & Evaluar y factorizar el jacobiano & Resolver con $B_k$ y actualizarla \\
Robustez global & Requiere salvaguardas & También requiere salvaguardas \\
Uso natural & Derivadas disponibles y costo moderado & Evaluar derivadas resulta costoso \\
\bottomrule
\end{tabular}
\normalsize
\end{frame}

\begin{frame}{Aplicación: equilibrio de Cournot}
Dos firmas eligen cantidades $q_1,q_2$ y $Q=q_1+q_2$. Considere
\[
P(Q)=Q^{-1/\alpha},
\qquad
C_i(q_i)=\frac12c_iq_i^2.
\]
El beneficio de la firma $i$ es
\[
\pi_i(q_1,q_2)=P(Q)q_i-\frac12c_iq_i^2.
\]
Las condiciones de primer orden producen el sistema
\[
f_i(q_1,q_2)=P(Q)+P'(Q)q_i-c_iq_i=0,
\qquad i=1,2.
\]
\end{frame}

\begin{frame}{Jacobiano del modelo de Cournot}
Como
\[
P'(Q)=-\frac1\alpha Q^{-1/\alpha-1},
\qquad
P''(Q)=\frac1\alpha\left(\frac1\alpha+1\right)Q^{-1/\alpha-2},
\]
el jacobiano tiene entradas
\[
\begin{aligned}
\frac{\partial f_i}{\partial q_i}
&=2P'(Q)+P''(Q)q_i-c_i,\\
\frac{\partial f_i}{\partial q_j}
&=P'(Q)+P''(Q)q_i,
\qquad j\neq i.
\end{aligned}
\]
Newton usa estas derivadas. Broyden puede comenzar con este jacobiano y
actualizarlo durante el cálculo.
\end{frame}

\begin{frame}{Restricciones y validación económica}
Una solución numérica debe comprobar más que un residuo pequeño.
\begin{itemize}
  \item Factibilidad: $q_1,q_2\geq0$ y $Q>0$.
  \item Equilibrio: $\norm{F(q^*)}$ debe ser pequeño.
  \item Estabilidad numérica: el resultado debe resistir cambios razonables en $q_0$.
  \item Interpretación: precio, cantidades y beneficios deben permanecer en escalas plausibles.
\end{itemize}
Una búsqueda amortiguada puede impedir pasos que produzcan cantidades negativas.
\end{frame}

\begin{frame}{Actividad de clase}
Para
\[
F(x,y)=\mat{x^2+y^2-1\\x-y},
\qquad (x_0,y_0)=(1,0.5),
\]
realice lo siguiente:
\begin{enumerate}
  \item construya el jacobiano;
  \item calcule dos pasos de Newton;
  \item use $B_0=J_F(x_0)$ y calcule una actualización de Broyden;
  \item compare $\norm{F(x_k)}_2$ en ambos métodos;
  \item explique por qué el método converge a la solución positiva.
\end{enumerate}
\end{frame}

\begin{frame}{Ideas centrales}
\begin{itemize}
  \item Newton convierte cada iteración no lineal en un sistema lineal.
  \item El jacobiano describe el acoplamiento local entre variables y ecuaciones.
  \item El amortiguamiento controla pasos poco fiables lejos de la solución.
  \item Broyden impone una condición secante mediante una actualización de rango uno.
  \item Residuo, tamaño del paso, condicionamiento y factibilidad completan la validación.
\end{itemize}
\end{frame}

\end{document}
